\section{Method Theory: DSR + DSRM Applied}
\label{sec:dsr-dsrm}

Design Science Research is the paradigm this thesis works in. The DSR foundation, together with the two alternatives it is compared against (positivist and interpretive research), is set out in \Cref{sec:dsr}. The short version is that neither alternative on its own fits this project. A purely positivist study would need to measure overtime and idle-time effects across many Norwegian transport companies in production deployment, which the project's stage does not support. A purely interpretive study would describe coordinator practice in rich detail without ever producing the runnable platform Admmit's task required. DSR is the third paradigm that takes care of both: an artefact plus the knowledge it embodies. RP is precisely such an artefact.

The methodology used in this thesis is \textcite{peffers2007dsrm}'s Design Science Research Methodology (DSRM), which structures DSR as six activities: problem identification and motivation, objectives for a solution, design and development, demonstration, evaluation, and communication \parencite[p.~46]{peffers2007dsrm}. Peffers note that the activities are nominally sequential rather than strictly so; researchers are not expected to walk in order from activity 1 through activity 6 \parencite[p.~56]{peffers2007dsrm}.

\begin{description}
  \item[Problem identification and motivation.] The seven coordinator interviews surfaced a recurring pattern: overtime was handled reactively, idle time between assignments was invisible in the tools coordinators used, and load balancing depended on the coordinator's memory rather than any structured overview. Admmit's task description framed the same problem from the owner side, asking for an algorithm-assisted planning platform that could make utilisation visible and adjustable. \Cref{sec:interview-findings} describes this in detail; \Cref{sec:efficiency} discusses it under the Efficiency anchor.
  \item[Objectives for a solution.] The objectives are stated under the three locked anchor concepts. \textbf{Efficiency} is realised by the artefact reducing overtime, reducing idle time between assignments, and reducing uneven load between drivers. \textbf{Trust/control} is realised as the coordinator's authority to inspect, modify, accept, or reject any algorithm-generated assignment. \textbf{Adaptability} is delivered through configurable constraint weights and per-company assignment rules within a single deployed instance that can serve companies with different operational rules.
  \item[Design and development.] The artefact is an algorithm-assisted planning platform built on Next.js (a React-based web framework), tRPC (a typed remote-procedure-call layer between the client and the server), and PostgreSQL (a relational database). Its solver layer integrates three engines, each invoked as a subprocess: a greedy constructive heuristic, the OR-Tools CP-SAT solver (Google's open-source constraint-programming engine), and Timefold (an open-source metaheuristic constraint-solver platform). These three are referred to as the three engines hereafter, and \Cref{sec:iter-multiengine} details them. A drag-and-drop timeline lets the coordinator inspect, modify, accept, or reject each assignment. The platform was built across eight named iterations, each labelled by origin (Admmit-mandated, interview-driven, or designer / domain-knowledge), described in \Cref{sec:iterations}.
  \item[Demonstration.] Demonstration takes the form of a multi-engine benchmark on synthetic datasets. The datasets cover the realistic fleet ranges of the seven interviewed companies, from small fleets of around ten drivers up to multi-department operations of around forty-five vehicles. The benchmark shows how the three engines behave when given the same problem instances. The dataset design and benchmark protocol are set out in \Cref{sec:evaluation-framework}.
  \item[Evaluation.] What the artefact is validated against is how solver approaches compare under realistic constraint combinations. It is not validated against whether the visibility gap is real, or whether the HITL stance is necessary. The benchmark validates solver performance independently of those two whether-questions. The requirements traceability matrix records, for each functional and non-functional requirement, whether it was implemented and how its behaviour was checked. What is not validated in this thesis is documented in \Cref{sec:limitations}: the artefact is not deployed in production, real-world utilisation gains are not measured, and the override flow is not user-tested with coordinators. The validation-versus-evaluation distinction that bounds these claims is set out under DSR theory in \Cref{sec:dsr} and applied to the benchmark in \Cref{sec:evaluation-framework}.
  \item[Communication.] The research is communicated through two channels. The first is this thesis, addressed to the academic audience at NTNU and presenting RP as a concrete instance of a DSR artefact alongside the design knowledge embedded in its multi-engine and HITL choices. The second is a written recommendation to Admmit covering deployment readiness, including the requirements not yet validated through user testing (\Cref{sec:limitations}) and the per-company configurability through which Adaptability is delivered.
\end{description}

The order in which the six activities were carried out did not follow Peffers' nominal sequence. Peffers note four possible entry points into the DSRM cycle and explicitly identify objective-centered initiation, in which an industry or research need fixes key objectives before activity 1 is complete, as a legitimate starting point \parencite[p.~56]{peffers2007dsrm}. RP fits that pattern: Admmit's task fixed HITL and the multi-tenant architecture (the artefact elements that realise Trust/control and Adaptability) before activity 1 was complete, and the eight iterations of activity 3 ran in parallel with later refinements of the objectives.


